\section{Integrasi API dan Menghubungkan Front-End dengan Back-End (request/response, data terstruktur) (Sub-CPMK 3.3)}

Integrasi front-end dengan back-end melalui API menggunakan \textbf{Fetch API} atau \textbf{XMLHttpRequest} di JavaScript. Fetch API adalah cara modern untuk mengirim HTTP request dari browser ke server tanpa me-reload halaman (\textit{AJAX}). Data dikirim dan diterima dalam format JSON. Pendekatan ini memisahkan tampilan (front-end) dari logika bisnis (back-end), memungkinkan pengembangan independen dan mendukung arsitektur \textit{Single Page Application} (SPA) \cite{phpManual}.

\begin{webcode}[language=JavaScript, caption={JavaScript Fetch API untuk mengakses REST API}]
// GET - Ambil daftar produk
async function ambilProduk() {
  const response = await fetch('/api/produk');
  const data = await response.json();
  data.forEach(p => {
    console.log(`${p.nama}: Rp ${p.harga}`);
  });
}

// POST - Tambah produk baru
async function tambahProduk(nama, harga) {
  const response = await fetch('/api/produk', {
    method: 'POST',
    headers: { 'Content-Type': 'application/json' },
    body: JSON.stringify({ nama, harga })
  });
  const result = await response.json();
  console.log('Created ID:', result.id);
}

// DELETE - Hapus produk
async function hapusProduk(id) {
  await fetch(`/api/produk/${id}`, { method: 'DELETE' });
  console.log('Produk dihapus');
}

// Panggil
ambilProduk();
tambahProduk('Keyboard', 250000);
\end{webcode}

